\begin{abstract}
Agentic LLMs for electronic design automation (EDA) are already deployed in-house, with a frontier large language model acting as the \emph{orchestrator}---the agent that faces the user, understands intent, decomposes the work, and invokes subagents.
% 中文：面向电子设计自动化（EDA）的智能体 LLM 已在公司内部署，由一个大模型担任 \emph{编排者}（orchestrator）——直面用户、理解意图、拆解工作并调用子智能体。
Cost, not feasibility, is now the binding constraint: every EDA subtask routed to the orchestrator consumes tokens and serving budget that scale with design-flow volume.
% 中文：当前约束是成本而非可行性：每个路由到前沿 编排者 的 EDA 子任务都消耗随设计流程规模增长的 token 与 serving 预算。
Privacy compounds the constraint: netlists, GDSII layouts, and design rules are trade secrets that must not leave the premises, so the on-site subtask is served by a small local model or not at all.
% 中文：隐私进一步加剧了这一约束：网表、GDSII 版图与设计规则是商业机密，不得离开本地，因此本地子任务只能由小型本地模型来服务，否则便无从谈起。
We present \textbf{BaiZe}, a $\approx$2.2B hybrid state-space agentic model built from scratch to serve as a small, local subagent beneath the orchestrator: it translates the orchestrator's intent into native EDA scripts (PyAether, SKILL, Tcl) and operates each tool.
% 中文：我们提出 \textbf{BaiZe}——一个约 22 亿参数的混合状态空间智能体模型，从零构建，作为编排者之下的一个小型本地子智能体：它把编排者的意图翻译成原生 EDA 脚本（PyAether、SKILL、Tcl）并操作每一个工具。
Its 56 layers interleave Mamba-2 state-space and MLP layers with four attention layers, so only the latter carry a session-growing state; decoding is linear-time and per-subtask serving cost is governed by session memory rather than prompt throughput.
% 中文：其 56 层将 Mamba-2 状态空间层与 MLP 层同 4 个注意力层交错排列，因此只有后者携带随会话增长的状态；解码为线性时间，每子任务的 serving 成本由会话内存而非 prompt 吞吐决定。
A matched from-scratch comparison against a dense transformer selects this architecture (11.6\% fewer parameters, lower loss at a matched 1k-step budget, and $+10.5\%$ prefill throughput), and a budget-constrained search fixes the pre-training recipe, now scaled to a 2.94T-token production run.
% 中文：通过与稠密 Transformer 的、条件对齐的从零训练对比选定了该架构（参数少 11.6\%、在同等 1k 步预算下损失更低、prefill 吞吐高 10.5\%），并由受预算约束的搜索确定预训练配方，现已放大到 2.94T token 的生产训练。
The backbone now stands at training loss 2.151 after 30{,}660 iterations ($\approx$598K tok/s).
% 中文：骨干现已训练至第 30,660 迭代、训练损失 2.151（约 598K tok/s）。
On a 449-task EDA-scripting benchmark spanning three languages and four commercial runtimes~\cite{liu2026zhulong}, we test whether BaiZe can offload subtasks from the deployed orchestrators at comparable task accuracy and materially lower serving cost.
% 中文：在一个覆盖三种语言、四种商业 runtime 的 449 任务 EDA 脚本基准上，我们检验 BaiZe 能否以可比的任务精度、显著更低的 serving 成本，从已部署的 orchestrator 卸载子任务。
Later stages are specified as protocols with pre-specified criteria; their measurements are in progress, and we release the full recipe and evaluation protocol.
% 中文：后续阶段以协议形式给出，并配有事先写定的评估准则；其测量正在进行中，我们公开完整配方与评估协议。
\end{abstract}

\begin{IEEEkeywords}
small language model, agentic LLM, hybrid state-space architecture, EDA subagent, cost-efficient inference
\end{IEEEkeywords}